\lesson{13}{Mar 16 2022 Wed (10:40:43)}{Find all possible rational zeros: zeros theorem}{Unit 4}

\begin{theorem}[Rational Zero Theorem]
  \label{thrm:rational_zero_theorem}

  Let $f(x) = a_{n} x^{n} + a_{n - 1} x^{n - 1} + \ldots + a_1 x + a_0$ be a
  polynomial with integer coefficients and $a_{n} \ne 0$. If $\frac{p}{q}$ is a
  zero of $f(x)$ is written in lowest terms, then $p$ is a factor of $a_0$, and
  $q$ is a factor of $a_n$ 

  In other words, the only possibilities for rational zeros of $f(x)$ are
  numbers of the form $\frac{p}{q}$ (in lowest terms), where $p$ is a factor of
  $a_0$, and $q$ is a factor of $a_n$
\end{theorem}

\begin{example}
  \label{exm:rational_zero_theorem}

  Let's list all of the possible rational zeros for the following polynomial:
  $f(x) = x^{3} + 2x^{2} - 5x - 6$

  \begin{equation}
    \label{eqt:rational_zero_theorem}
    \begin{split}
      \frac{p}{q} &= \frac{\pm 1, \pm 2, \pm 3, \pm 6}{\pm 1} \\
                  &= \pm 1, \pm 2, \pm 3, \pm 6 \\
    \end{split}
  .\end{equation}

  Here, $p$ is factors of the constant term, which in this case is  $-6$ and $q$
  is the leading coefficient.

  Now, let's start to check:

  \begin{equation}
    \label{eqt:checking_zeros_from_rational_zero_theorem}
    \begin{split}
      f(1) &= (1)^{3} + 2(1)^{2} - 5(1) - 6 \\
           &= 1 + 2 - 5 - 6 \\
           &= 3 - 11 = -8 \\
    \end{split}
  .\end{equation}

  Notice that $1$ doesn't give you $0$. Therefore, $1$ isn't a zero of the
  polynomial $f(x)$.

  \begin{equation}
    \label{eqt:checking_zeros_from_rational_zero_theorem}
    \begin{split}
      f(1) &= (2)^{3} + 2(2)^{2} - 5(2) - 6 \\
           &= 8 + 8 - 10 - 6 \\
           &= 16 - 16 = 0 \\
    \end{split}
  .\end{equation}

  Now, instead of just brute forcing to check all of the zeros, let's use one of
  the zeros we found, which was $2$. For this, we have to use synthetic
  division (\ref{sub_sub_sec:synthetic_division}):

  \polyhornerscheme[x=2]{x^3+2x^2-5x-6}

  That gives us:

  \begin{equation}
    \label{eqt:after_synthetic_division}
    \begin{split}
      x^{2} + 4x + 3 &= 0 \\
      (x + 3)(x + 1) &= 0 \\
      x + 3 = 0 \qquad&\qquad x + 1 = 0 \\
      x &= -3, -1
    \end{split}
  .\end{equation}

  So, here are our final zeros: $x = 2, -3, -1$.
\end{example}

\newpage
